%% HAND-WRITTEN fragment -- NOT generated by maestro2tex (raw Spectre Monte
%% Carlo, not a Maestro results database; there is no maestro2tex config for any
%% *_eis_* fragment).  Precedent: fig_eis_error.tex, fig_pixtop_noise.tex.
%%
%% Drop-in source: ~/chips/castalia/simruns/eisfigs2_20260826/build/
%% fig2_mc_calibration.tex, standalone preamble and in-figure stat blocks
%% removed -- the statistics are in the caption here.
%% Data: data/eis_mc_{law,dc}.dat, copied from eisfigs2_20260826/data/
%% fig2_hist_{law,dc}.dat, binned there from
%% ~/chips/castalia/simruns/eisdeep_20260826/mc_report/typ_n0_vs_cal11k.csv
%% (per-sample errBlaw_0.1 and errBdc_0.1, 649 rows); sigma and its exact chi2
%% 95 % interval from mc_report/typ_n0_vs_cal11k.txt.  The decomposition quoted
%% in the caption is mc_report/decomp_proc.txt (n = 500), decomp_mism.txt
%% (n = 500) and cal_n0.txt (n = 654).  Netlist provenance: mc/typ_n0/input.scs,
%% section castalia_pm_25.  Bin widths 2.5 % (left) and 2e-4 % (right).
%%
%% TWO INDEPENDENT X-SCALES, ON PURPOSE.  Not a shared axis and not a twin
%% axis: the two distributions differ in width by 1.4e4, so one axis would draw
%% the right panel as a single rule.  The caption says so, because a reader who
%% compares the two spreads by eye reads the figure backwards.
%%
%% Key numbers.  Stored law sigma = 10.242 % [9.714 - 10.832], mean -7.627 %,
%% range -32.06 to +30.63 %, n = 649.  Per-chip dc calibration sigma =
%% 0.00065 % [0.00061 - 0.00068], mean +0.00325 %, range +0.00166 to
%% +0.00623 %, n = 649 -- a factor of 1.58e4, and indistinguishable from
%% architecture A on the same run.  Underlying R_f 19 727 ohm +- 2195 =
%% 11.129 % [10.555 - 11.770] over the same 649 dies.  Decomposition, on the
%% cal_n0 cell and on R_f itself rather than on the impedance error: process
%% only 10.907 % (n = 500), mismatch only 0.4262 % (n = 500), quadrature
%% 10.916 % against 11.119 % [10.547 - 11.757] measured with both (n = 654) --
%% inside the interval, so the published sigma(process) > sigma(both) inversion
%% was sampling error.  DO NOT read 10.24 % and 10.91 % as one number measured
%% twice: different quantities on different cells.
%%
%% Caveats that had to survive into the caption: process AND mismatch,
%% castalia_pm_25, NOMINAL CORNER ONLY, 37 C, seed 12345, donominal=yes;
%% castalia_pm_25 holds tt_mim, so the 22 pF feedback capacitor does not vary
%% and this bounds R_f and not the pole; 0.1 Hz is BELOW the feedback pole,
%% above which a dc calibration stops working and the ratiometric architecture
%% is what rescues it; no samples censored or held out, all 649 plotted; the
%% range is this run's own (mc/typ_n0) and not a union over the campaign.
%% Library cell names are kept out of the rendered prose.
\providecommand{\MaestroRoot}{include/analog/}
%% mtxE and mtxC are the chapter's own palette slots and carry the same
%% values here; \providecolor is therefore a no-op if an earlier fragment
%% already defined them, which is what is wanted.  Any colour that is NOT a
%% chapter slot must be given a name of its own -- see fig_eis_settling.tex.
\providecolor{mtxE}{HTML}{E87BA4}
\providecolor{mtxC}{HTML}{1BAF7A}
\begin{figure}[htbp]
  \centering
  {\pgfplotsset{compat=1.18}%
  \pgfplotsset{
    mcax/.style={
      width=0.45\linewidth, height=4.3cm,
      ymajorgrids,
      major grid style={line width=0.2pt, draw=black!14},
      axis line style={line width=0.4pt, draw=black!45},
      tick style={line width=0.4pt, draw=black!45},
      tick align=outside,
      label style={font=\small}, tick label style={font=\small},
      title style={font=\small},
      ybar interval=0.9,
      xlabel={\(|Z|\) magnitude error~[\si{\percent}]},
      ylabel={Dies per bin}, ymin=0,
      scaled x ticks=false,
    }}%
  \begin{tikzpicture}
    \begin{axis}[mcax,
      xmin=-35, xmax=35, ymax=84, xtick={-30,-20,-10,0,10,20,30},
      title={\textbf{Stored design law}},
    ]
      \addplot[fill=mtxE, draw=white, line width=0.4pt]
        table[x=x, y=n] {\MaestroRoot data/eis_mc_law.dat};
    \end{axis}
  \end{tikzpicture}\hfill
  \begin{tikzpicture}
    \begin{axis}[mcax,
      xmin=0.0015, xmax=0.0065, ymax=109,
      xtick={0.002,0.003,0.004,0.005,0.006},
      xticklabel style={/pgf/number format/fixed,
                        /pgf/number format/precision=3},
      title={\textbf{One dc calibration per chip}},
    ]
      \addplot[fill=mtxC, draw=white, line width=0.4pt]
        table[x=x, y=n] {\MaestroRoot data/eis_mc_dc.dat};
    \end{axis}
  \end{tikzpicture}}
  \caption[{Impedance error over 649 dies under two calibration
  policies.}]{Impedance-magnitude error over \num{649} Monte Carlo dies at
  \SI{0.1}{\hertz} on the typical electrode at gain code 0, with process and
  mismatch together at the nominal corner and \SI{37}{\celsius}, under two
  calibration policies: a stored \(\SI{18}{\kilo\ohm} + \SI{6}{\kilo\ohm}\,N\)
  design law (left) and one dc measurement of the feedback resistance per chip
  (right). \(\sigma\) falls from \SI{10.24}{\percent}
  (\SI{95}{\percent} interval \numrange{9.71}{10.83}) to
  \SI{0.00065}{\percent} (\numrange{0.00061}{0.00068}), a factor of
  \num{1.6e4}, and the right panel is indistinguishable from measuring the
  current at every frequency --- the bound, which needs a current reference the
  die does not carry. \textbf{The two \(x\)-axes are independent and are not to
  be compared by eye}: the left spans \SI{70}{\percent} and the right
  \SI{0.005}{\percent}, and both are the same \num{649} dies, load, gain code
  and run. The plotted range is that run's own, \(-32.1\,\%\) to
  \(+30.6\,\%\), and not a union over the campaign; no sample is censored. The
  statistical model section holds the capacitor family at typical, so the
  \SI{22}{\pico\farad} feedback capacitor does not vary here and the figure
  bounds the feedback resistance rather than the pole --- which matters because
  \SI{0.1}{\hertz} is below that pole, and above it a dc calibration stops
  working at all (Section~\ref{sss:eis-accuracy}).}
  \label{fig:eis-mccal}
\end{figure}
